\documentclass[11pt,a4paper]{article}

\usepackage[T1]{fontenc}
\usepackage{lmodern}
\input{glyphtounicode}
\pdfgentounicode=1
\usepackage{amsmath,amssymb,amsthm,mathtools}
\usepackage{booktabs}
\usepackage{doi}
\usepackage{enumitem}
\usepackage{geometry}
\usepackage{hyperref}
\usepackage{microtype}
\usepackage{xcolor}

\geometry{margin=2.5cm}
\hypersetup{
  colorlinks=true,
  linkcolor=blue,
  citecolor=blue,
  urlcolor=blue,
  pdftitle={A Finite Spectral Architecture for Mass-Square Splitting and Mixing},
  pdfauthor={Jerome Beau}
}

\newtheorem{theorem}{Theorem}[section]
\newtheorem{proposition}[theorem]{Proposition}
\newtheorem{lemma}[theorem]{Lemma}
\newtheorem{corollary}[theorem]{Corollary}
\theoremstyle{definition}
\newtheorem{definition}[theorem]{Definition}
\theoremstyle{remark}
\newtheorem{remark}[theorem]{Remark}
\newtheorem{example}[theorem]{Example}

\newcommand{\C}{\mathbb{C}}
\newcommand{\F}{\mathbb{F}}
\newcommand{\Fq}{\mathbb{F}_q}
\newcommand{\G}{\mathrm{SL}(2,\Fq)}
\newcommand{\Heis}{\mathrm{Heis}_3(\Fq)}
\newcommand{\End}{\mathrm{End}}
\newcommand{\Ind}{\mathrm{Ind}}
\newcommand{\Span}{\mathrm{Span}}
\newcommand{\Stab}{\mathrm{Stab}}
\newcommand{\tr}{\mathrm{tr}}
\newcommand{\Id}{\mathbf{1}}
\newcommand{\Dplus}{\mathcal{D}_{+}}
\newcommand{\Dminus}{\mathcal{D}_{-}}
\newcommand{\AK}{A_K}
\newcommand{\HK}{H_c(K)}

\title{\textbf{A Finite Spectral Architecture for Mass-Square Splitting and Mixing:}\\
\textbf{Exact Doublets and Protected Sectors from Weil Admissibility}\\[0.5em]
\large O33 --- Spectral Admissibility Programme}

\author{J\'er\^ome Beau\\[0.3em]
\small Independent Researcher\\
\small \texttt{jerome.beau@cosmochrony.org}}
\date{2026}

\begin{document}

\maketitle

\begin{abstract}
We give an exact arithmetic and representation-theoretic description of a finite Weil
filtration generated by three-character Heisenberg fingerprints.
For every odd prime $q$ and non-zero total character $c=c_1+c_2+c_3$, the cumulative Fourier
mode set at depth $n$ is the cyclic sumset
\[
  K_n=c[-(n+1),n+1]+\{0,\pm(c_2+c_3)\}+\{0,\pm c_3\}.
\]
This closed form turns the filtration into a union of at most nine cyclic intervals.
It yields an exact spacing criterion for boundary porosity and reduces the full Weil
stabiliser to the multiplicative symmetry group of $K_n$.
Weyl quantisation then identifies the induced Mackey carrier internally as a direct summand
of $\End(\C^q)$.
Its generic principal-series constituents occur in canonical multiplicity-two pairs with
commutant $M_2(\C)$.
The finite Harper operator selects one axis in every such doublet and produces the universal
normalised split $1\pm q^{-1/2}$.
The eliminated-sector residue supplies a second axis precisely when an explicit cyclotomic
phase determinant is non-zero.
The odd Weyl-symbol sector is itself a pure doublet reservoir for $q\equiv1\pmod4$, but it is
exactly protected by the parity-even algebra generated by the deposited projectors and
Harper dynamics.
An orbit-divisibility sieve constrains every multiplicative stabiliser, and uniform
constructions give exact $C_4$ and $C_6$ exceptional families, together with a third,
near-saturation $C_4$ family with four-point complement.
An unconditional almost-periodicity lemma shows that every multiplicative stabiliser shifts
the mode set with additive defect exactly twice the run count, and a run-correspondence
argument with a Diophantine descent proves that the stabiliser is $\{\pm1\}$ whenever every
run and gap of the mode set is longer than the run count.
Completeness of the exceptional orders is thereby reduced to a degenerate boundary regime;
exhaustive finite audits find no order beyond six there through $q=73$ and also expose exact
vanishings of selected boundary character sums.

\noindent\textbf{Interpretive status.}
The construction supplies finite two-level spectral splitting, mixing frames, and
mass-square kinematics at an operator stratum.
It does not derive Dirac masses, an absolute mass scale, Standard-Model generations,
hypercharge, or a matter-sector Higgs mechanism.
\end{abstract}

\noindent\textbf{Keywords.}
Finite Weil representation; Heisenberg group; Mackey induction; Hecke algebra;
principal series; cyclic sumsets; spectral filtration; Harper operator; mixing;
parity protection.

\newpage
\tableofcontents

\newpage
\section{Problem, Results, and Epistemic Scope}
\label{sec:introduction}

The finite Weil representation is simultaneously a model of exact harmonic analysis over
finite fields and the carrier used by the O12 spectral-admissibility
computation~\cite{Weil1964,Gerardin1977,Beau2026a16}.
The present paper isolates a mathematical structure that was hidden by the original
Gram--Schmidt implementation: the cumulative O12 filtration is a Fourier coordinate
filtration whose mode sets have a closed additive form.

Three questions then become finite and exact.
First, which subgroup of $\G$ stabilises a proper filtration space?
Second, where can multiplicity-two factors occur once the irreducible parity sectors of the
finite Weil module itself have scalar commutant?
Third, which deposited operators select axes inside those factors?

The answers form the following chain.
The mode space has a complete stabiliser controlled by a multiplicative set stabiliser.
That set stabiliser is rigid: an additive almost-periodicity argument proves it equal to
$\{\pm1\}$ outside a degenerate boundary regime.
The Weyl-symbol algebra embeds the relevant induced representation into
$\End(\C^q)$.
Mackey induction and the Weyl pairing of principal series create canonical $M_2(\C)$
commutants at this operator stratum~\cite{Mackey1952,Iwahori1964}.
Harper dynamics gives a first axis, while boundary porosity can give a second.
Parity simultaneously protects the complementary odd sector.

The paper separates four epistemic levels.

\begin{enumerate}[label=(\roman*)]
  \item The sumset formula, spacing theorem, stabiliser formula, additive rigidity theorem,
  Weyl-symbol decomposition, and parity closure are proved uniformly.
  \item The Hecke decompositions are standard finite-group consequences, with conventions
  stated explicitly.
  \item Exact cyclotomic claims are made either uniformly from their formulas or for the
  cases actually certified.
  \item Multi-prime frequencies and physical readings are labelled computational or
  interpretive and are not promoted to theorems.
\end{enumerate}

No identification with a Standard-Model representation is made.
In particular, a multiplicity doublet in an operator algebra is not by itself a weak-isospin
doublet of fermionic matter.

\section{Finite Heisenberg and Weil Setup}
\label{sec:setup}

Let $q$ be an odd prime and write elements of $\Heis$ as $(a,b,z)$ with product
\begin{equation}
  (a,b,z)(a',b',z')=(a+a',b+b',z+z'+ab').
  \label{eq:heisenberg-product}
\end{equation}
Let $X=(1,0,0)$ and $Y=(0,1,0)$, and use the symmetric Cayley generating set
$\{X^{\pm1},Y^{\pm1}\}$.

Fix a non-zero central character $c\in\Fq^*$ and set
$\psi_c(t)=\exp(2\pi i ct/q)$.
The Schr\"odinger representation on $V_c=\C^{\Fq}$ is taken in the convention
\begin{equation}
  \bigl(\rho_c(a,b,z)f\bigr)(x)=\psi_c\!\bigl(z+b(x-a)\bigr)f(x-a).
  \label{eq:schrodinger}
\end{equation}
Its extension to $\G$ is the finite Weil representation.
Only conjugation identities are used below, so the scalar metaplectic normalisation does not
affect the statements.

For $v=(a,b)\in\Fq^2$, define the symmetrically ordered displacement
\begin{equation}
  \bigl(W_c(a,b)f\bigr)(x)
  =\psi_c\!\left(b\left(x-\frac{a}{2}\right)\right)f(x-a).
  \label{eq:weyl-displacement}
\end{equation}
This is the finite-field analogue of Weyl's symmetric quantisation
prescription~\cite{Weyl1927}.
These $q^2$ operators are an orthogonal basis of $\End(V_c)$ and satisfy
\begin{equation}
  \rho_c(g)W_c(v)\rho_c(g)^{-1}=W_c\bigl(\varphi(g)v\bigr),
  \qquad g\in\G.
  \label{eq:weyl-covariance}
\end{equation}
In the coordinates of~\eqref{eq:weyl-displacement}, the linear phase-space action is
\begin{equation}
  \varphi(g)=R_0g^{-\mathsf{T}}R_0,
  \qquad R_0=\operatorname{diag}(1,-1).
  \label{eq:phase-space-action}
\end{equation}

For a triple $(c_1,c_2,c_3)\in(\Fq^*)^3$, the O12 fingerprint vectors are products of
three Schr\"odinger vectors and hence single Fourier characters.
We assume throughout that
\begin{equation}
  c:=c_1+c_2+c_3\ne0.
  \label{eq:total-character}
\end{equation}
The cumulative filtration space is therefore
\begin{equation}
  W_{<n}=\Span\{\chi_k:k\in K_n\},
  \qquad \chi_k(x)=q^{-1/2}\psi_1(kx),
  \label{eq:fourier-filtration}
\end{equation}
for a finite set $K_n\subset\Fq$.

\section{Closed Sumset Form and the Spacing Theorem}
\label{sec:sumset}

Set
\begin{equation}
  u=c_2+c_3,
  \qquad v=c_3.
  \label{eq:u-v}
\end{equation}

\begin{theorem}[Closed mode-set formula]
  \label{thm:closed-mode-set}
  At cumulative BFS depth $n\geq0$,
  \begin{equation}
    K_n=c[-(n+1),n+1]+\{0,\pm u\}+\{0,\pm v\}
    \quad\text{in }\Fq.
    \label{eq:closed-mode-set}
  \end{equation}
  In particular, $K_n=-K_n$.
\end{theorem}

\begin{proof}
  The second Heisenberg coordinate is additive:
  $b(gs)=b(g)+b(s)$.
  On the symmetric generating set, $b(s)$ belongs to $\{0,\pm1\}$.
  The BFS ball of depth $n$ realises every integer $b(g)\in[-n,n]$, because powers of
  $Y$ realise the endpoints and intermediate values.
  Conversely, each generator changes the second coordinate by at most one, so every ball
  element satisfies $|b(g)|\leq n$.

  Write $t_i=b(s_i)$ for the three post-shell generator choices.
  The three successive second coordinates are
  \[
    b_1=b(g)+t_1,
    \quad b_2=b(g)+t_1+t_2,
    \quad b_3=b(g)+t_1+t_2+t_3.
  \]
  Hence the fingerprint frequency is
  \[
    c_1b_1+c_2b_2+c_3b_3
    =c\bigl(b(g)+t_1\bigr)+ut_2+vt_3.
  \]
  As $b(g)$ and $t_1$ vary, $b(g)+t_1$ fills $[-(n+1),n+1]$.
  The variables $t_2,t_3$ independently fill $\{0,\pm1\}$.
  This proves~\eqref{eq:closed-mode-set}.
  Each summand is centrally symmetric, proving $K_n=-K_n$.
\end{proof}

Rescale by $c^{-1}$ and define
\begin{equation}
  \widetilde u=u/c,
  \quad \widetilde v=v/c,
  \quad
  S=\{a\widetilde u+b\widetilde v:a,b\in\{-1,0,1\}\}.
  \label{eq:centre-set}
\end{equation}
The rescaled set $c^{-1}K_n$ is a union of cyclic intervals of radius $n+1$ centred on the
at most nine distinct points of $S$.

Order the points of $S$ cyclically and denote their positive cyclic spacings by
$\delta_1,\ldots,\delta_{|S|}$.
Define the rescaled flanked-hole set
\begin{equation}
  H(K_n)=\{h\notin c^{-1}K_n:h-1,h+1\in c^{-1}K_n\}.
  \label{eq:flanked-holes-rescaled}
\end{equation}

\begin{theorem}[Spacing-parity criterion]
  \label{thm:spacing}
  Before saturation,
  \begin{equation}
    H(K_n)\ne\varnothing
    \quad\Longleftrightarrow\quad
    2n+4\in\{\delta_1,\ldots,\delta_{|S|}\}.
    \label{eq:spacing-criterion}
  \end{equation}
  Consequently an odd spacing never produces a flanked hole.
  An even spacing $\delta$ produces one at the unique depth
  $n=(\delta-4)/2$, provided that depth is pre-saturation.
\end{theorem}

\begin{proof}
  Consecutive windows of radius $n+1$ whose centres have cyclic distance $\delta_i$ leave
  exactly $\delta_i-2(n+1)-1=\delta_i-(2n+3)$ uncovered sites.
  A flanked hole is precisely a complementary component of cardinality one.
  This gives $\delta_i-(2n+3)=1$, which is~\eqref{eq:spacing-criterion}.
\end{proof}

\section{Complete Weil Stabiliser}
\label{sec:stabiliser}

For a non-empty proper set $K\subset\Fq$, define
\begin{equation}
  A_K=\{a\in\Fq^*:aK=K\}.
  \label{eq:AK}
\end{equation}

\begin{theorem}[Complete stabiliser]
  \label{thm:stabiliser}
  The stabiliser in the finite Weil action of the Fourier coordinate space
  $W_K=\Span\{\chi_k:k\in K\}$ is
  \begin{equation}
    \Stab_{\G}(W_K)
    =\left\{
      \begin{pmatrix}a&0\\t&a^{-1}\end{pmatrix}
      :t\in\Fq,\ a\in A_K
    \right\}.
    \label{eq:complete-stabiliser}
  \end{equation}
  It has order $q|A_K|$.
\end{theorem}

\begin{proof}
  In the Fourier basis, the lower unipotent radical acts diagonally and the diagonal torus
  sends the label $k$ to a non-zero scalar multiple of $k$.
  Hence a lower-triangular element preserves $W_K$ exactly when its multiplier lies in
  $A_K$.

  If $g$ is outside the lower Borel, its Weil kernel is a non-degenerate quadratic Gauss
  kernel.
  Every matrix coefficient of the image of a Fourier basis vector is then non-zero.
  Thus the image has full Fourier support and cannot lie in a non-empty proper coordinate
  subspace.
  The two statements give~\eqref{eq:complete-stabiliser}.
\end{proof}

Because Theorem~\ref{thm:closed-mode-set} gives $K_n=-K_n$, one always has
$\{\pm1\}\subseteq A_{K_n}$.
For the eight deposited cases, $A_{K_n}=\{\pm1\}$ and the stabiliser is the abelian group
$\langle-I\rangle U_-$ of order $2q$.

This equality is generic in the computational audit but not universal.
For example, an exact depth-zero case at $q=29$ has $A_K\cong C_4$, while one at $q=31$
has $A_K\cong C_6$.
The general theorem is therefore~\eqref{eq:complete-stabiliser}, not the special value
$A_K=\{\pm1\}$.

The exceptional groups nevertheless obey a sharp elementary sieve.

\begin{proposition}[Orbit-divisibility sieve]
  \label{prop:AK-divisibility}
  For every proper $K_n$, the group $A_{K_n}$ is cyclic of even order and
  \begin{equation}
    |A_{K_n}|\mid\gcd(q-1,|K_n|-1).
    \label{eq:AK-divisibility}
  \end{equation}
  Equivalently, $|A_{K_n}|$ divides both $|K_n|-1$ and $q-|K_n|$.
  In particular, if the greatest common divisor in~\eqref{eq:AK-divisibility} is two,
  then $A_{K_n}=\{\pm1\}$.
\end{proposition}

\begin{proof}
  The multiplicative group $\Fq^*$ is cyclic, so every subgroup is cyclic.
  Central symmetry gives $-1\in A_{K_n}$, hence the order is even.
  The origin belongs to $K_n$ and is fixed, while the action of $A_{K_n}$ on
  $K_n\setminus\{0\}$ is free.
  Thus $|A_{K_n}|$ divides $|K_n|-1$.
  The same free action on the complement gives divisibility of $q-|K_n|$.
\end{proof}

There are two uniform exceptional families.
They explain the orders four and six without asserting that these are the only possible
orders for the full equal-radius nine-centre class.

\begin{proposition}[Uniform $C_4$ and $C_6$ families]
  \label{prop:exceptional-families}
  Suppose first that $a^2=-1$ in $\Fq$ and $q>5$.
  At depth zero, the normalised parameters $\widetilde u=0$ and
  $\widetilde v=a$ give a proper set with $A_K\cong C_4$.

  Suppose instead that $a^2+a+1=0$, $a\ne1$, and $q>19$.
  At depth zero, the parameters $\widetilde u=a$ and
  $\widetilde v=a+1$ give a proper set with $A_K\cong C_6$.
\end{proposition}

\begin{proof}
  In the first family,
  \[
    K=\{m+na:m,n\in\{-1,0,1\}\}.
  \]
  Multiplication by $a$ sends $(m,n)$ to $(-n,m)$, so $C_4\subseteq A_K$.
  The nine elements are distinct for $q>5$.
  Proposition~\ref{prop:AK-divisibility} leaves only orders four and eight.
  If the order were eight, then $K\setminus\{0\}$ would be the order-eight subgroup and
  would contain $1+a$.
  But $(1+a)^8=16$, so this would require $q\mid15$, contrary to $q>5$.

  In the second family, write an element as $m+na$.
  The set $K$ consists of the origin and three six-element orbits under
  \begin{equation}
    (a+1)(m+na)=(m-n)+ma.
    \label{eq:C6-action}
  \end{equation}
  The nineteen elements are distinct for $q>19$: the norm of every non-zero difference
  is one of $1,3,4,7,9,12,13,16$.
  Hence $C_6\subseteq A_K$, and Proposition~\ref{prop:AK-divisibility} leaves only orders
  six and eighteen.
  Order eighteen would make $K\setminus\{0\}$ a multiplicative subgroup.
  Since $2\in K$, it would force $4\in K$.
  Substitution in the nineteen forms reduces this possibility to prime divisors of
  $7,12,13,19,21,28,31$.
  Together with the cubic-root condition, this leaves
  $q\in\{7,13,19,31\}$; the first three cases are saturated.
  At $q=31$, closure would further force $8\in K$, while direct substitution gives
  possible divisors only from $43,52,57,67,73,84,91$, so $8\notin K$.
  Thus the proper stabiliser has order six.
\end{proof}

Section~\ref{sec:rigidity} proves that such exceptions are confined to a degenerate boundary
regime, and adds a third uniform family at the near-saturation end of that regime.

\section{Additive Rigidity of the Multiplicative Stabiliser}
\label{sec:rigidity}

Throughout this section $K\subset\Fq$ is proper and non-empty, with maximal cyclic intervals
(\emph{runs}) $I_1,\ldots,I_m$ and complementary maximal intervals (\emph{gaps})
$G_1,\ldots,G_m$, listed in cyclic order.
Write $I_i=[x_i,y_i]$ and $L_i=|I_i|$.
For the mode sets of Theorem~\ref{thm:closed-mode-set} one has $m\leq9$.

\begin{lemma}[Unconditional almost-periodicity]
  \label{lem:almost-period}
  For every $a\in A_K$,
  \begin{equation}
    |K\cap(K+a)|=|K|-m,
    \qquad
    |K\,\triangle\,(K+a)|=2m.
    \label{eq:almost-period}
  \end{equation}
\end{lemma}

\begin{proof}
  Let $C(d)=|K\cap(K-d)|$.
  For $a\in A_K$ the bijection $x\mapsto a^{-1}x$ of $\Fq$ maps $K\cap(K-ad)$ onto
  $K\cap(K-d)$, so $C(ad)=C(d)$ for every $d$.
  Each run of length $L_i$ contains exactly $L_i-1$ elements $x$ with $x+1\in K$, so
  $C(1)=|K|-m$.
  Taking $d=1$ gives the first equality, and the second follows from $|K+a|=|K|$.
\end{proof}

Every multiplicative stabiliser of a nine-centre mode set is therefore an additive
almost-period of $K$ with symmetric-difference defect at most eighteen.
The multiplicative question becomes a question about the additive boundary of $K$.

\begin{lemma}[Run correspondence]
  \label{lem:run-correspondence}
  Assume every run and every gap of $K$ has length at least $m+1$.
  Then for every $a\in A_K$ there are a permutation $\sigma$ of $\{1,\ldots,m\}$ and
  non-negative integers $\lambda_i,\rho_i,\alpha_i,\beta_i$ with
  $\sum_i(\lambda_i+\rho_i)=\sum_i(\alpha_i+\beta_i)=m$, such that each
  $M_i=[x_i+a+\lambda_i,\,y_i+a-\rho_i]$ is non-empty, satisfies
  $M_i=(I_i+a)\cap K\subseteq I_{\sigma(i)}$, and starts at $x_{\sigma(i)}+\alpha_i$.
  Consequently
  \begin{equation}
    ma\equiv t\pmod q,
    \qquad
    t=\sum_i\alpha_i-\sum_i\lambda_i,
    \qquad
    0<|t|\leq m.
    \label{eq:residue-bound}
  \end{equation}
\end{lemma}

\begin{proof}
  Translation by $a$ preserves the run--gap pattern, so the sets $J_i=I_i+a$ are the runs of
  $K+a$, separated by gaps of the same lengths as those of $K$.
  Let $U_i=J_i\setminus K$; Lemma~\ref{lem:almost-period} gives
  $\sum_i|U_i|=|(K+a)\setminus K|=m$.

  A gap $G$ of $K$ meeting $J_i$ either contains an endpoint of $J_i$, or satisfies
  $G\subseteq J_i$ or $J_i\subseteq G$.
  The last two options put at least $\min(|G|,|J_i|)\geq m+1>m$ points into $U_i$, which is
  impossible.
  Distinct gaps are disjoint, so at most one gap contains each endpoint of $J_i$.
  Hence $U_i$ is the union of a prefix of length $\lambda_i\geq0$ and a suffix of length
  $\rho_i\geq0$ of $J_i$, and $M_i=J_i\setminus U_i=J_i\cap K$ is an interval of length
  $L_i-\lambda_i-\rho_i\geq L_i-m\geq1$.
  Being a non-empty interval inside $K$, $M_i$ lies in a unique run $I_{\sigma(i)}$.

  Suppose $\sigma(i)=\sigma(j)=k$ with $i\ne j$, and choose $M_i$ and $M_j$ adjacent among
  the middles contained in $I_k$, with $M_j$ following $M_i$ inside $I_k$.
  If $\rho_i>0$, the point following $M_i$ lies in $U_i$, hence outside $K$; then $M_i$ ends
  at the right endpoint of $I_k$ and no point of $I_k$ lies beyond it, contradicting
  $M_j\subseteq I_k$.
  So $\rho_i=0$ and $M_i$ ends where $J_i$ ends.
  The gap of $K+a$ that follows $J_i$ has length at least $m+1$ and lies between two points
  of the interval $I_k$, hence inside $I_k\subseteq K$.
  Its points are outside $K+a$, so $|K\setminus(K+a)|\geq m+1$, contradicting
  Lemma~\ref{lem:almost-period}.
  Hence $\sigma$ is injective, therefore a permutation.

  If $p\in I_k\cap(K+a)$, then $p\in J_l\cap K=M_l\subseteq I_{\sigma(l)}$ for some $l$, so
  $l=\sigma^{-1}(k)$: every run satisfies $I_k\cap(K+a)=M_{\sigma^{-1}(k)}$.
  Writing $\alpha_i\geq0$ and $\beta_i\geq0$ for the numbers of points of $I_{\sigma(i)}$
  before and after $M_i$, the identity $|K\setminus(K+a)|=m$ gives
  $\sum_i(\alpha_i+\beta_i)=m$.

  The left endpoint of $M_i$ equals both $x_i+a+\lambda_i$ and $x_{\sigma(i)}+\alpha_i$.
  Summing over $i$ and using $\sum_ix_{\sigma(i)}=\sum_ix_i$ gives
  $ma\equiv\sum_i(\alpha_i-\lambda_i)\pmod q$.
  The bound $|t|\leq m$ follows from $0\leq\sum_i\alpha_i\leq m$ and
  $0\leq\sum_i\lambda_i\leq m$.
  Finally $t\ne0$: the hypotheses force $q\geq2m(m+1)>m$, so $ma\ne0$ in $\Fq$.
\end{proof}

\begin{theorem}[Regular-regime rigidity]
  \label{thm:rigidity}
  If every run and every gap of $K$ has length at least $m+1$, then
  $A_K\subseteq\{\pm1\}$.
  For the centrally symmetric mode sets of Theorem~\ref{thm:closed-mode-set} this gives
  $A_K=\{\pm1\}$.
\end{theorem}

\begin{proof}
  The hypotheses force $q\geq2m(m+1)>2m^2$, so the Diophantine threshold is automatic.
  Fix $a\in A_K$ of order $h$ and apply Lemma~\ref{lem:run-correspondence} to each power
  $a^j\in A_K$: there are integers $t_j$ with $ma^j\equiv t_j\pmod q$ and
  $0<|t_j|\leq m$.
  Multiplying two such congruences gives
  $t_it_j\equiv m^2a^{i+j}\equiv mt_{i+j}\pmod q$; both sides lie in
  $[-m^2,m^2]\subset(-q/2,q/2)$, so $t_it_j=mt_{i+j}$ holds in $\mathbb{Z}$.
  By induction $t_j=t_1^j/m^{j-1}$ for all $j\geq1$.
  Since $a^h=1$, the residue $t_h\equiv m\pmod q$ with $|t_h|\leq m$ forces $t_h=m$, so
  $(t_1/m)^h=1$.
  The only rational roots of unity are $\pm1$, hence $t_1=\pm m$ and
  $a\equiv m^{-1}t_1=\pm1$.
\end{proof}

\begin{corollary}[Reduction to the degenerate regime]
  \label{cor:degenerate-reduction}
  Every proper mode set with $A_K\ne\{\pm1\}$ has a run or a gap of length at most
  $m\leq9$.
  The completeness question for the exceptional orders is thereby reduced to this degenerate
  boundary regime.
\end{corollary}

All known exceptional families indeed live there.
The depth-zero $C_4$ and $C_6$ families of Proposition~\ref{prop:exceptional-families} have
runs of length three.
At the opposite, near-saturation end of the regime sits a third uniform family with a
four-point complement.

\begin{proposition}[Near-saturation $C_4$ family]
  \label{prop:near-saturation}
  Let $K^{\mathrm{c}}=\{w,w+1,-w-1,-w\}$ with the four elements distinct and non-zero.
  Then $A_K\ne\{\pm1\}$ if and only if
  \begin{equation}
    2w^2+2w+1\equiv0\pmod q,
    \label{eq:near-saturation}
  \end{equation}
  in which case $A_K=\langle(w+1)/w\rangle\cong C_4$ and $K^{\mathrm{c}}$ is a single free
  $A_K$-orbit.
  Equation~\eqref{eq:near-saturation} is solvable exactly when $q\equiv1\pmod4$.
\end{proposition}

\begin{proof}
  One has $A_K=A_{K^{\mathrm{c}}}$, and $0\in K$ makes the action of $A_K$ on
  $K^{\mathrm{c}}$ free: $az=z$ with $z\ne0$ forces $a=1$.
  Hence $|A_K|$ divides four.
  If $|A_K|=4$, its generator $a$ spans the unique order-four subgroup of $\Fq^*$, so
  $a^2=-1$, and the orbit of $w$ is all of $K^{\mathrm{c}}$.
  From $aw\ne-w$ (else $a=-1$) we get $aw=\pm(w+1)$, hence $a=\pm(w+1)/w$, and $a^2=-1$ is
  equivalent to $(w+1)^2\equiv-w^2$, which is~\eqref{eq:near-saturation}.
  Conversely, if $w$ solves~\eqref{eq:near-saturation}, set $a=(w+1)/w$.
  Then $a^2=-1$ and
  $aw=w+1$, $a(w+1)=a^2w=-w$, $a(-w)=-w-1$, $a(-w-1)=w$,
  so $a\in A_{K^{\mathrm{c}}}=A_K$, and the divisibility bound gives $A_K\cong C_4$.
  Solvability: the discriminant of $2w^2+2w+1$ is $-4$, a square modulo $q$ exactly when
  $q\equiv1\pmod4$.
\end{proof}

The family is realised by deposited-class mode sets.
At $q=73$, the normalised block $(\widetilde u,\widetilde v)=(46,3)$ reaches, at depth
$n=8$, the mode set with $K^{\mathrm{c}}=\{13,14,59,60\}$ and $w=13$; analogous
realisations occur at $q=109$ (depth $11$) and $q=173$ (depth $10$).
The sieve of Proposition~\ref{prop:AK-divisibility} is consistent: $|K|=q-4$ gives
$\gcd(q-1,|K|-1)=\gcd(q-1,4)$.

\begin{proposition}[Small complements]
  \label{prop:small-complement}
  If $|K^{\mathrm{c}}|=4$, then $|A_K|$ divides four; if $A_K=\langle a\rangle$ has order
  four, then $a^2=-1$ and $K^{\mathrm{c}}=\{z,az,-z,-az\}$ for any
  $z\in K^{\mathrm{c}}$.
\end{proposition}

\begin{proof}
  The action of $A_K$ on $K^{\mathrm{c}}$ is free as in
  Proposition~\ref{prop:near-saturation}, so $|A_K|$ divides $|K^{\mathrm{c}}|=4$.
  In the order-four case any point orbit has size four, hence equals $K^{\mathrm{c}}$.
\end{proof}

The audited exceptional cases with $|K^{\mathrm{c}}|=4$ all realise this orbit in the
adjacent form of Proposition~\ref{prop:near-saturation}.
Whether the sumset geometry excludes the remaining centrally symmetric gap profiles --- a
symmetric pair of singleton gaps together with one two-point gap, or four singleton gaps ---
is not settled here.
The adjacency statement is verified on the audited cases and is not promoted to a theorem.

\section{The Internal Mackey Carrier}
\label{sec:carrier}

Equation~\eqref{eq:weyl-covariance} makes the Weyl basis a permutation basis for the
$\G$-action by conjugation.
There are two phase-space orbits: the origin and the punctured plane.
The stabiliser of a non-zero vector is a unipotent subgroup $U$ of order $q$.

\begin{proposition}[Internal induced representation]
  \label{prop:internal-carrier}
  As a $\G$-module under conjugation,
  \begin{equation}
    \End(V_c)\cong\Id\oplus\Ind_U^{\G}(\Id).
    \label{eq:internal-induced}
  \end{equation}
  Splitting the punctured orbit by $v\mapsto-v$ gives
  \begin{equation}
    \End(V_c)\cong\Id\oplus\Dplus\oplus\Dminus,
    \qquad
    \dim\Dplus=\dim\Dminus=\frac{q^2-1}{2}.
    \label{eq:parity-decomposition}
  \end{equation}
\end{proposition}

\begin{proof}
  Map the delta function at $v\in\Fq^2$ to $W_c(v)$.
  Weyl orthogonality makes this a vector-space isomorphism
  $\C[\Fq^2]\to\End(V_c)$, and covariance makes it $\G$-equivariant.
  The origin gives the trivial summand.
  The punctured plane is $\G/U$, proving~\eqref{eq:internal-induced}.
  The central element $-I$ acts freely on the punctured plane, so its two eigenspaces have
  equal dimension and give~\eqref{eq:parity-decomposition}.
\end{proof}

This proposition resolves a canonicity question.
The induced carrier is not an external representation added to the finite model.
It is already present in the bosonic operator algebra $\End(V_c)$.

\section{Hecke Doublets and Canonical Multiplicity Algebras}
\label{sec:hecke}

Let $Z=\{\pm I\}$ and let $H=ZU$.
The even sector is $\Dplus\cong\Ind_H^{\G}(\Id)$, while the odd sector is induced from the
non-trivial character of $Z$.
Induction in stages through a Borel $B=TU$ decomposes these modules into principal-series
representations.

For a character $\theta$ of $T/Z$, the Weyl element identifies the principal series
$\mathrm{PS}(\theta)$ and $\mathrm{PS}(\theta^{-1})$.
When $\theta\ne\theta^{-1}$, the two occurrences therefore form a canonical multiplicity
space $\C^2$.
The diagonal torus and the Weyl intertwiner generate $M_2(\C)$ on that space.
This is the relevant finite Hecke algebra, not merely the group algebra of the exchange
involution~\cite{Iwahori1964}.

\begin{proposition}[Even-sector commutant]
  \label{prop:even-commutant}
  If $q\equiv1\pmod4$, then
  \begin{equation}
    \End_{\G}(\Dplus)
    \cong\C^4\oplus M_2(\C)^{(q-5)/4}.
    \label{eq:even-commutant-one}
  \end{equation}
  If $q\equiv3\pmod4$, then
  \begin{equation}
    \End_{\G}(\Dplus)
    \cong\C^2\oplus M_2(\C)^{(q-3)/4}.
    \label{eq:even-commutant-three}
  \end{equation}
  In both cases the commutant dimension is $q-1$.
\end{proposition}

\begin{proof}
  Decompose $\Ind_H^{\G}(\Id)$ by the characters of $T/Z$.
  Generic characters occur in inverse pairs and contribute one $M_2(\C)$ block per pair.
  The self-inverse trivial and quadratic characters split according to their central sign.
  For $q\equiv1\pmod4$, both belong to the even sector and contribute four scalar
  constituents in total.
  For $q\equiv3\pmod4$, only the trivial branch contributes to the even sector, giving two
  scalar constituents.
  Counting the remaining inverse pairs gives~\eqref{eq:even-commutant-one} and
  \eqref{eq:even-commutant-three}.
\end{proof}

The original finite Weil module $V_c$ itself has only its two irreducible parity factors and
a two-dimensional scalar commutant~\cite{Korinman2013}.
The multiplicity-two structure therefore appears strictly at the operator stratum
$\End(V_c)$, consistently with the irreducibility of the two parity factors.

\section{Harper Orientation of the Doublets}
\label{sec:harper}

Define the Hermitian finite Harper operator~\cite{Harper1955}
\begin{equation}
  T_c=W_c(X)+W_c(-X)+W_c(Y)+W_c(-Y).
  \label{eq:harper}
\end{equation}
Its Weyl symbol is supported on two transverse Lagrangian lines.
The two same-line contributions are orthogonal, while the transverse contribution is a
finite Gauss sum.

Let $\theta$ be a generic even torus character and let
\begin{equation}
  \tau_{\theta,c}
  =\sum_{x\in\Fq^*}\theta^{-1}(x)\psi_c(x).
  \label{eq:gauss-sum}
\end{equation}
The classical Gauss identity gives $|\tau_{\theta,c}|^2=q$; its general finite-field
setting is part of the Davenport--Hasse theory~\cite{DavenportHasse1935}.

\begin{proposition}[Universal Harper split]
  \label{prop:harper-split}
  On every generic multiplicity doublet of $\Dplus$, the reduced Harper density has equal
  diagonal entries and non-zero off-diagonal entry proportional to
  $\tau_{\theta,c}$.
  After division by its scalar part, its two eigenvalues are
  \begin{equation}
    1-\frac{1}{\sqrt q},
    \qquad
    1+\frac{1}{\sqrt q}.
    \label{eq:harper-split}
  \end{equation}
\end{proposition}

\begin{proof}
  Homogeneous projection onto the pair
  $\{\theta,\theta^{-1}\}$ makes the same-line terms orthogonal with equal norms.
  Additive orthogonality reduces the two transverse terms to
  $2(q-1)\tau_{\theta,c}/q$.
  The diagonal term is $2(q-1)$ in the same normalisation.
  Since $|\tau_{\theta,c}|=\sqrt q$, the ratio of the off-diagonal modulus to the diagonal
  term is $q^{-1/2}$, proving~\eqref{eq:harper-split}.
\end{proof}

The phase of $\tau_{\theta,c}$ fixes the intertwiner gauge in each factor.
After this canonical Gauss normalisation, the two eigenlines are the symmetric and
antisymmetric Weyl combinations.
The splitting scale tends to zero as $q\to\infty$; no continuum mass gap follows without a
separate scaling prescription.

\section{Boundary Residues and Factorwise Mixing}
\label{sec:boundary}

Let $P_K$ be the Fourier-coordinate projector onto $W_K$ and define
\begin{equation}
  E_K=-P_KT_c(1-P_K)T_cP_K.
  \label{eq:residue}
\end{equation}
This is a negative semidefinite Hermitian square.
In the Fourier basis, the clock part of $T_c$ commutes with $P_K$ and drops out of
$(1-P_K)T_cP_K$.
Only the shifts by $\pm c$ remain.

Define the unscaled flanked-hole set
\begin{equation}
  H_c(K)=\{h\notin K:h-c,h+c\in K\}.
  \label{eq:holes}
\end{equation}
The diagonal part of $E_K$ is the negative outgoing shift degree.
Its Weyl symbol is real, even, and supported on one Lagrangian line, so it is scalar in every
generic multiplicity factor.
The off-diagonal part joins the two endpoints around each $h\in H_c(K)$.

For a generic even character $\theta$, set
\begin{equation}
  C_\theta(K)=\sum_{h\in H_c(K)}\theta\!\left(\frac{h}{2c}\right)
  \label{eq:boundary-character-sum}
\end{equation}
and
\begin{equation}
  \Delta_{\theta,c}(K)
  =C_\theta(K)\overline{\tau_{\theta,c}}
  -\overline{C_\theta(K)}\tau_{\theta,c}.
  \label{eq:phase-determinant}
\end{equation}

\begin{theorem}[Exact factorwise mixing criterion]
  \label{thm:mixing}
  The traceless multiplicity densities of $T_c$ and $E_K$ generate $M_2(\C)$ in the generic
  $\theta$ factor if and only if
  \begin{equation}
    \Delta_{\theta,c}(K)\ne0.
    \label{eq:mixing-criterion}
  \end{equation}
  If $H_c(K)=\varnothing$ or $C_\theta(K)=0$, then the second axis is absent.
\end{theorem}

\begin{proof}
  The boundary reduction above makes the off-diagonal component of the reduced $E_K$
  density proportional to $C_\theta(K)$.
  Proposition~\ref{prop:harper-split} makes the Harper off-diagonal component proportional
  to $\tau_{\theta,c}$.
  Two Hermitian $2\times2$ matrices generate $M_2(\C)$ exactly when their traceless
  three-component coefficient vectors are not collinear.
  Their oriented area is proportional to~\eqref{eq:phase-determinant}, giving the
  criterion.
\end{proof}

When~\eqref{eq:mixing-criterion} holds, orthogonalising the traceless part of $E_K$ against
that of $T_c$ and orienting the third axis by $i[T_c,E_K]$ gives a canonical ordered
$\mathfrak{su}(2)$ frame.
This construction is conjugation-covariant and therefore independent of the phase convention
used to realise the principal-series factor.

Theorem~\ref{thm:spacing} decides when a boundary coefficient can occur.
It does not by itself decide~\eqref{eq:mixing-criterion}: exact cyclotomic cancellation can
make $C_\theta(K)$ vanish, and a non-zero $C_\theta(K)$ must still pass the phase-determinant
test.

\section{Arithmetic Exceptions and Computational Certification}
\label{sec:audit}

The exact formulas were audited in three layers.
First, direct BFS enumeration at $q=29$ agrees with
Theorem~\ref{thm:closed-mode-set} for three deposited blocks and depths zero through two.
All eight previously deposited cases reproduce their exact dimensions and porous/smooth
dichotomy.

Second, a seeded sweep used 39 primes from 29 through 211 together with 307, with 30
admissible character blocks per prime.
It contains 1,170 blocks and 19,153 pre-saturation depth cases.
The spacing criterion has no failure.
An exhaustive audit over all normalised admissible parameter pairs at
$q=13,17,29,31,37,41,43,53$ also has no spacing failure.
These checks guard the proof but do not replace it.

The multiplicative stabiliser is $\{\pm1\}$ in the eight deposited cases.
The random sweep contains seven order-four exceptions.
An exhaustive audit of every normalised admissible pair and every pre-saturation depth for
all primes $13\leq q\leq73$ finds only orders two, four, and six.
The order-four cases occur for
$q=13,17,29,37,41,53,61,73$, while the order-six cases occur for
$q=31,37,43,61,67,73$.
A separate exact seeded audit through $q=499$ tested 12,234 cases surviving the divisor
sieve and found no order beyond six.
Proposition~\ref{prop:exceptional-families} was also checked for every applicable prime
through $997$.
Together with Theorem~\ref{thm:rigidity}, these finite audits reduce the conjecture that
every proper mode set of the present form has $A_K\cong C_2$, $C_4$, or $C_6$ to the
degenerate regime of Corollary~\ref{cor:degenerate-reduction}, where it remains open.

The rigidity chain of Section~\ref{sec:rigidity} was validated by the same exact machinery.
An exhaustive sweep over all normalised centre parameters and pre-saturation depths for
$q\leq73$ (256,229 cases) and a seeded sweep of sixty blocks per prime for
$74\leq q\leq250$ (42,878 cases) checked 599,020 instances of
Lemma~\ref{lem:almost-period}, 408,558 residue bounds of
Lemma~\ref{lem:run-correspondence}, and 204,279 instances of Theorem~\ref{thm:rigidity},
with no failure.
All 355 exceptional cases in these sweeps (307 of order four, 48 of order six) have a run
or a gap of length at most $m$, as Corollary~\ref{cor:degenerate-reduction} requires.

Third, every porous sweep case was checked in the exact cyclotomic ring
$\mathbb{Z}[\zeta_{q-1}]$ for zeros of~\eqref{eq:boundary-character-sum}.
Among 1,186 porous cases, exactly 50 contain a zero: 84 inverse-paired generic factors in
total.
Twenty-nine cases reduce to a two-term relation
$\theta(x_1/x_2)=-1$, and twenty-one to a three-term cyclotomic cancellation.

For the remaining 36,285 generic factors, a numerical phase scan finds no candidate
alignment in~\eqref{eq:phase-determinant}; the smallest relative determinant is about
$7.12\times10^{-6}$.
This last statement is numerical only.
Exact non-alignment has been certified for the eight deposited cases, not uniformly over the
full sweep.

In the random sample, 294 of 1,170 blocks have no porous pre-saturation depth.
The fraction $25\%$ is descriptive of the seeded sample and is not asserted as an
asymptotic density.

\section{Parity Protection and Calibrated Outlook}
\label{sec:conclusion}

Let $R$ be Schr\"odinger parity, $(Rf)(x)=f(-x)$.
Then
\begin{equation}
  RW_c(v)R^{-1}=W_c(-v).
  \label{eq:parity-weyl}
\end{equation}
The $R$-even and $R$-odd operator spaces are respectively
$\Id\oplus\Dplus$ and $\Dminus$.

\begin{theorem}[Parity-closure no-go]
  \label{thm:parity-closure}
  The deposited operators $P_{K_n}$, $T_c$, $(1-P_{K_n})T_cP_{K_n}$, and $E_{K_n}$ are
  $R$-even.
  The unital star-algebra they generate, including ordered compositions, functional calculus,
  and norm limits, has zero projection onto $\Dminus$.
\end{theorem}

\begin{proof}
  Theorem~\ref{thm:closed-mode-set} gives $K_n=-K_n$, hence $P_{K_n}$ is even.
  The four terms in~\eqref{eq:harper} occur in opposite pairs, hence $T_c$ is even.
  Products and sums of even operators are even because
  $R(AB)R^{-1}=(RAR^{-1})(RBR^{-1})$.
  Adjoints, functional calculus, and norm limits preserve the same identity.
  The odd projection therefore vanishes throughout the generated algebra.
\end{proof}

For $q\equiv1\pmod4$, the odd-sector commutant is
\begin{equation}
  \End_{\G}(\Dminus)\cong M_2(\C)^{(q-1)/4}.
  \label{eq:odd-commutant}
\end{equation}
For $q\equiv3\pmod4$, the quadratic character moves to the odd sector and contributes its
two scalar constituents, giving
\begin{equation}
  \End_{\G}(\Dminus)\cong\C^2\oplus M_2(\C)^{(q-3)/4}.
  \label{eq:odd-commutant-three}
\end{equation}
Thus the protected sector is not empty.
It is a pure family of multiplicity doublets when $q\equiv1\pmod4$, while the
$q\equiv3\pmod4$ branch also contains two scalar constituents.
An algebraically natural odd operator would be
\begin{equation}
  T_c^{\mathrm{odd}}
  =i\bigl(W_c(X)-W_c(-X)+W_c(Y)-W_c(-Y)\bigr).
  \label{eq:odd-harper}
\end{equation}
It requires a choice inside each generator pair.
The directed Q11 frontier explicitly selects no subset of generators, and its arrow is
radial in BFS depth~\cite{Beau2026q11of}.
Projective temporal ordering orders that same radial depth variable~\cite{Beau2026pto}.
Neither construction presently supplies the phase-space-odd lift
\eqref{eq:odd-harper}.

\subsection*{Interpretive outlook}

The proved mathematical structure is a parity-graded operator architecture.
The even sector carries a universal finite split and, when the boundary determinant is
non-zero, a second axis and a complete two-level frame.
The odd sector carries further doublets but is protected from all operators in the deposited
even algebra.

It is structurally tempting to compare the first sector with bosonic symmetry breaking and
the protected sector with a missing chiral datum.
That comparison is interpretive only.
The construction does not identify electroweak charges, does not transport the doublets to a
fermionic matter carrier, and does not produce an unsquared left--right mass operator.
The residue $E_K$ remains of $-M^*M$ type.

The exact advance toward masses is consequently limited but concrete: the theory derives
two-level mass-square kinematics, a finite split of size $q^{-1/2}$, and arithmetic mixing
frames indexed by the filtration.
An absolute scale, a continuum law, and a chiral descent remain open.

The stabiliser problem is now localised.
Theorem~\ref{thm:rigidity} proves $A_K=\{\pm1\}$ whenever every run and gap of the mode set
is longer than the run count, so exceptional orders occur only in the degenerate boundary
regime of Corollary~\ref{cor:degenerate-reduction}, which hosts the depth-zero $C_4$ and
$C_6$ families and the near-saturation $C_4$ family of
Proposition~\ref{prop:near-saturation}.
The remaining completeness question is to classify that regime: prove or refute that the
exceptional orders four and six exhaust it.
The other principal questions are a uniform arithmetic criterion for
$\Delta_{\theta,c}(K)\ne0$ and the asymptotic distribution of porous depths.

\bibliographystyle{plain}
\bibliography{cosmochrony-bibliography}

\end{document}
